% =====================================================================
%  Выводы по лабораторной работе.
% =====================================================================

\labpart{Выводы}

В ходе выполнения лабораторной работы был освоен механизм контроля
доступа на уровне приложений и систем управления базами данных -- в
отличие от первой части, где права ограничивались правами файловой
системы, здесь для MongoDB реализована и экспериментально подтверждена
трёхуровневая политика безопасности «администратор -- пользователь --
гость», выраженная собственными, специфичными для этой СУБД средствами:
встроенными и пользовательскими ролями с привилегиями на уровне базы
данных и коллекции. Практически показано, что в основе этого механизма
лежит принцип наименьших привилегий: пользователю и гостю явно и по
отдельности выдаётся ровно тот объём прав, который необходим для их
роли, а не урезается из более широкого набора.

Семантически лабораторная работа подтвердила, что задача администратора
не сводится к однократной установке ПО: то, что MongoDB по умолчанию
разрешает подключение вовсе без аутентификации, означает, что первым
шагом всегда является явное включение и настройка контроля доступа --
то же наблюдение, что и в первой части работы, только теперь на уровне
не операционной системы, а конкретного приложения. Использование
Docker-контейнера для установки также оказалось практичным решением:
оно позволило безопасно выполнять реальные, а не имитационные команды
установки и настройки, и гарантировало полное и воспроизводимое удаление
всех созданных объектов на завершающем этапе.
